\begin{table}[htbp]
\centering
\caption{Summary Statistics}
\label{tab:summary}
\begin{tabular}{lcc}
\hline\hline
Variable & Mean & SD \\
\hline
\multicolumn{3}{l}{\textit{Panel A: Trade outcomes}} \\
Mercury imports (\$1,000s) & 103.2 & 443.6 \\
Gold exports (\$1,000s) & 658898.6 & 2205451.5 \\
\hline
\multicolumn{3}{l}{\textit{Panel B: Country characteristics}} \\
GDP per capita (\$) & 2529 & 3182 \\
Population (millions) & 21.3 & 31.1 \\
Rule of law & -0.72 & 0.64 \\
Control of corruption & -0.65 & 0.63 \\
\hline
\multicolumn{3}{l}{\textit{Panel C: Treatment variables}} \\
ASGM country (\%) & 42.6 & 49.5 \\
Minamata ratifier (\%) & 10.5 & 30.7 \\
EU mercury share (pre-ban) & 0.412 & 0.412 \\
\hline
Observations & \multicolumn{2}{c}{864} \\
Countries & \multicolumn{2}{c}{54} \\
Years & \multicolumn{2}{c}{16} \\
\hline\hline
\end{tabular}
\begin{tablenotes}[flushleft]\small
\item \textit{Notes:} Summary statistics computed over the full descriptive sample (54 African countries, 2005--2020, $N=864$). Regression samples differ: EU ban analysis uses 2005--2015 excluding 2011 ($N=540$); Minamata analysis uses 2005--2023 ($N=1{,}003$). Mercury and gold trade from UN Comtrade (HS 280540 and 7108). Governance indicators from World Governance Indicators. EU mercury share is the average share of mercury imports sourced from EU member states during 2005--2010.
\end{tablenotes}
\end{table}
